\chapter{From an accepted interpretation to verified Java behavior}
\label{ch:verification}

% BEGIN SOURCE UNIT M07:00
\section{The specification becomes the comparison target}
\label{merge:M07:00}


Once reviewers accept a particular interpretation, a different verification
problem begins. The source question was whether the rule expresses the relevant
regulatory requirement. The implementation question is whether the deployed
program preserves that accepted rule for the inputs and histories it supports.
Keeping these questions separate allows strong engineering evidence without
misrepresenting it as legal proof.

The accepted object is larger than a Boolean formula. It includes the source
versions, activity profile, factual definitions, subject identity, time policy,
unknown and conflict semantics, exception structure, result interpretation and
coverage limits. If any of these is omitted, two implementations can agree on a
formula while giving different advice in practice. A Java adapter that labels
all rebates as discounts has changed the rule's application even if its Boolean
operator code is flawless.

The verification target must therefore name the complete bundle and its
permitted input domain. The bundle is immutable and content-addressed. The
comparison includes the decision, missing and blocking facts, errors, source
references and execution trace. Different backends have different execution
identities, so equality of their entire serialized result hashes is not the
right criterion. Compare the specified semantic fields and independently verify
each backend's own result identity.

The existing workbench already provides a finite RuleIR interpreter, a Java
runtime, generated policy classes and a release workflow. The ensemble proposed
in Chapter~\ref{ch:ensemble} is an extension before that delivery path. It should
produce reviewed candidates that the existing validator can accept, rather than
silently expanding the runtime language whenever a model proposes a new legal
construct. Unsupported meaning remains an explicit engineering and review issue.
% END SOURCE UNIT M07:00

% BEGIN SOURCE UNIT P-03-product:02
\section{The intermediate representation}
\label{merge:P-03-product:02}


A JSON document becomes an executable specification only when its expressions
have defined meanings. RuleIR 0.1 fixes a finite, acyclic expression language and
a separate event profile. The following sections give its expression, evidence,
event, storage and service rules in the manuscript. The JSON schemas and
\path{docs/specs/v0.1/} files provide machine-checkable structures and detailed
operational copies of those contracts. They do not select an alternative logic
or leave missing-value behavior to the implementer.

Each rule needs more than its condition. It must identify the source document and
revision, paragraph or page, exact supporting span, language, source digest and
retrieval date. It must record the legal actor, activity, product, client class,
relevant time and interpretation status. It must distinguish a derived fact from
a duty, permission or prohibition. A duty also needs an actor, trigger, action,
timing rule and evidence by which performance can be assessed. A missing source
deadline remains missing; operational escalation times belong to a separately
identified bank policy.

Definitions should be reusable and versioned. A field for portfolio value must
refer to the adopted portfolio definition, relevant ownership rules, valuation
date and currency conversion convention. The application mapping records where
each fact comes from, its freshness and who can resolve a conflict. A rule whose
necessary definition remains unresolved can be reviewed and discussed, but it
cannot be released as an executable control for that use.

The complete financial-condition bundle is supplied in
\path{fixtures/spi-financial.bundle.json} beneath the specification directory.
It declares two money inputs, an unconditional scope for this isolated
subcondition, the expression, interpretation record and a verified span in
Annex 1 paragraph 3.1. The excerpt below is a complete expression node from
that bundle. Node identifiers make its trace reproducible.

\begin{minipage}{\textwidth}
\begin{lstlisting}[caption={The portfolio comparison in the supplied JSON bundle. Amounts are HK cents.}]
{"node_id":"portfolio.ge", "op":"compare", "cmp":"ge",
 "left":{"node_id":"portfolio.input", "op":"fact",
         "name":"portfolio"},
 "right":{"node_id":"portfolio.threshold", "op":"literal",
          "type":"money_hkd", "value":"4000000000"}}
\end{lstlisting}
\end{minipage}

The second comparison uses \texttt{net\_assets\_ex\_home} and
\texttt{"8000000000"}; an \texttt{any} node combines them. These inputs denote
amounts normalized under adopted wealth definitions, not any account balance
supplied under the same name. Release requires those definitions and applicability
to be reviewed. The narrow result name, \texttt{spi.financial}, identifies a
subcondition rather than a complete SPI assessment or trade authorization.
% END SOURCE UNIT P-03-product:02

% BEGIN SOURCE UNIT P-03a-rule-contract:00
\subsection{Types and the expressions an agent may generate}
\label{merge:P-03a-rule-contract:00}


The scalar types are Boolean, integer, HKD money and Gregorian date. Integer
and money values use canonical decimal strings; monetary units are cents.
There are no floating-point values, implicit casts or known nulls. A date must
be valid, including its leap-year constraints. Foreign-currency normalization
records the rate, valuation instant and rounding policy before producing an HKD
input. A missing exchange rate therefore leaves an unknown amount.

The expression grammar contains typed literals, fact and rule references,
\texttt{all}, \texttt{any}, \texttt{not}, equality/inclusive/strict comparison,
addition, subtraction, rational scaling, conditional choice and explicit defaults.
Addition and subtraction require matching integer or money types. Scaling uses
a nonnegative integer numerator and a positive denominator; a known result must
divide exactly. A half-cent allocation returns \texttt{E\_INEXACT\_SCALE} until
an external reviewed rounding policy resolves it. Date arithmetic is kept in
the separate scheduler so ordinary library behavior cannot choose a legal
anniversary convention unnoticed.

An implementer can express the grammar below directly as tagged records. Every
\texttt{Expr} additionally carries a unique \texttt{node\_id}. A rule is
\texttt{(id, type, scope, body, interpretation\_id, source\_span\_ids)};
\texttt{scope} is Boolean and \texttt{body} has the declared result type.
\begin{lstlisting}[caption={Complete operator grammar for RuleIR 0.1; repeated arguments are ordered. The Java demonstration implements a stated subset.}]
Type = bool | integer | money_hkd | date
Expr = literal(Type, value)
     | fact(name) | rule(name)
     | all([Expr, ...]) | any([Expr, ...]) | not(Expr)
     | compare(eq | ge | gt, left, right)
     | add(left, right) | sub(left, right)
     | scale(arg, numerator, denominator)
     | if(condition, then_value, else_value)
     | default(base, [Exception, ...])
Exception = (exception_id, guard, value,
             interpretation_id, source_span_ids)
\end{lstlisting}
A comparison requires two matching integer, money or date operands. Dates
compare in calendar order; Boolean comparisons are not admitted in this version.
A conditional requires a Boolean condition and matching branch types. Default
values match their base type, and all guards are Boolean. An empty exception
list simply selects the base. There are no loops, quantifiers, dynamic names or
arbitrary function calls. More elaborate assessment and aggregation enter
through separately reviewed fact mappings.

Every node has a unique identifier. References must resolve, types must agree
and the complete dependency graph must be acyclic. A rule reference may refer
only to an unconditional definition; a scoped result cannot quietly become a
Boolean false in another expression. Authors can expose the scope explicitly
as a named Boolean definition. Unknown operators and additional JSON properties
are rejected. This keeps a model from extending the language merely by emitting
a plausible field name.
% END SOURCE UNIT P-03a-rule-contract:00

% BEGIN SOURCE UNIT M07:01
\section{Facts are evidence-bearing values}
\label{merge:M07:01}


A production fact is not simply a field such as \texttt{gift=true}. It has a type,
value or uncertainty status, evidence references, validity interval and recording
time. The validity interval says when the evidence applies to the underlying
situation. The recording time says when the system knew it. These times answer
different questions and both are needed for historical reconstruction.

Suppose the bank learns on 10 October that a fee record effective on 1 October
was incorrect. A decision replay asking what the system knew on 5 October should
not use the later correction. A corrected assessment made on 10 October may use
it to reassess the earlier event. The two assessments should coexist. Replacing
the old record would erase the evidence needed to explain the original decision.

The existing normalization boundary selects records valid at \texttt{valid\_at}
and recorded no later than \texttt{known\_at}. A correction names the evidence it
supersedes. Two admissible, non-superseded records with different values create a
conflict. Arrival order is not a source-priority rule. If both records agree,
their evidence may be combined without inventing a conflict.

Intervals are half-open: $[a,b)$ includes the start and excludes the end. If a
record is valid until midnight, it is not valid at that midnight. This convention
avoids overlapping adjacent intervals when one record ends exactly as another
begins. Tests must exercise equality at both endpoints, not merely dates well
inside an interval. All assessment timestamps use the canonical UTC format with
six fractional digits; legal calendar dates are separate values.

Unknown facts have reasons. Missing evidence, stale evidence, an unreviewed
classification and unresolved event order can all lead to unknown, but they call
for different remedies. A host should not display the same uninformative message
for every case. Conflict similarly preserves the competing evidence identifiers.
The evaluator does not decide which source is legally superior without a
reviewed normalization rule.

This design puts factual interpretation where it belongs: at a visible boundary.
For the circular, classifying a benefit as a gift or fee discount is an evidenced
assessment. The decision engine consumes that assessment under its approved
definition. If the ensemble has not resolved the definition, the adapter cannot
supply a known value merely to satisfy a required input field.
% END SOURCE UNIT M07:01

% BEGIN SOURCE UNIT P-03a-rule-contract:01
\subsection{Evidence is an input to computation, with its own history}
\label{merge:P-03a-rule-contract:01}


A normalized fact is known, unknown or conflicting. Known facts have an exact
value, evidence identifiers, a validity interval and a recorded timestamp.
Unknown facts explain whether evidence is missing, stale or awaiting assessment.
Conflicts preserve the disagreeing records. Two admissible records with the
same value can combine their evidence; different values require resolution.
Arrival order is not a ranking of legal authority.

Every request supplies \texttt{valid\_at}, the instant being assessed, and
\texttt{known\_at}, the latest evidence the assessment may use. Intervals are
half-open: an amount valid until midnight is unavailable at that midnight.
A correction names the record it supersedes and becomes visible only when its
recorded time is eligible. Thus a corrected assessment can use new information
while the original decision remains reproducible.

Before Boolean evaluation, the engine checks conflicts in the static transitive
dependency closure, including branches that might later be skipped. A conflicting
portfolio amount therefore returns \texttt{CONFLICT} even if the net-asset route
could establish the financial condition. An unrelated conflict does not block
that rule. This conservative policy is a proposed engineering choice; it is
deliberately different from ordinary unknown-value propagation.
% END SOURCE UNIT P-03a-rule-contract:01

% BEGIN SOURCE UNIT M07:02
\section{Exact arithmetic prevents a second kind of ambiguity}
\label{merge:M07:02}


The existing scalar types are Boolean, integer, Hong Kong dollar money in cents,
and Gregorian date. Integer and money values use canonical decimal strings and
arbitrary-precision arithmetic. A monetary value of HK\$100 is represented as
\texttt{"10000"}. The string representation avoids a transport layer converting
the value through a floating-point number before the exact arithmetic begins.

Units are part of the contract. A number expressed in dollars and the same number
expressed in cents differ by a factor of one hundred. A compiler cannot infer
which was intended from the digits alone. The host mapping must identify the
source unit, conversion and evidence. Declaration-specific bounds can reject
implausible inputs, but a plausibility check is not a substitute for a reviewed
unit mapping.

Consider proportional ownership of a monetary amount. If the accepted operation
is exact scaling by a rational fraction $n/d$, the runtime computes the product
and quotient using integers. Let $x$ be cents. It checks
\begin{equation}
 nx=dq+r,\qquad 0\leq r<d
 \label{eq:exact-scale}
\end{equation}
for nonnegative $nx$; for signed values the implementation's quotient and
remainder convention must still make zero remainder the exactness test. A
nonzero remainder means that the result is not an integer number of cents. The
current contract returns \texttt{E\_INEXACT\_SCALE}; it does not choose a rounding
policy on the user's behalf.

If a bank requires rounding, that is a substantive extension. The rule must state
which quantity is rounded, at what stage, in which direction and under whose
policy. Rounding each component before summation can differ from rounding the
sum. Two correct implementations of different policies can therefore disagree.
The ensemble should recognize this as a specification issue rather than treating
one result as a numerical accident.

Foreign-currency conversion is outside RuleIR 0.1. A reviewed adapter supplies an
HKD amount with rate, valuation instant and rounding evidence. A missing rate
creates unknown evidence. It must not trigger a guessed conversion or reuse an
unboundedly old rate. The bank's data owner is responsible for the freshness and
valuation policy, while engineering verifies that the adapter implements it.

Dates require similar precision. A date is not a timestamp, and an obligation
measured in business days needs a calendar policy. The current expression
language compares validated dates but leaves date arithmetic to reviewed event
scheduling. This restricted design prevents a generic arithmetic library from
silently deciding whether a public holiday counts toward a regulatory deadline.
The appropriate calendar must be versioned with the assessment.
% END SOURCE UNIT M07:02

% BEGIN SOURCE UNIT M07:03
\section{Unknown, conflict and scope have an order}
\label{merge:M07:03}


The three-valued Boolean rules in Chapter~\ref{ch:proof} govern unknown values,
but they are not the entire runtime semantics. Before evaluating the requested
rule, the existing implementation checks conflicting evidence in its static
transitive dependency closure. That closure includes scope, body, exceptions and
branches that might later be skipped. A conflict in the closure returns conflict
even if another route could determine the Boolean result.

This is a conservative evidence policy. For example, an actor may appear outside
the selected scope, but a conflicting fact in the rule's static closure still
produces conflict before scope evaluation. The second-circular case that exercises
this behavior is deliberately different from ordinary short-circuit intuition.
The specification must teach it because an engineer using conventional Java
\texttt{\&\&} or \texttt{||} could otherwise implement a different result.

After conflict checking, scope is evaluated. False scope returns
\texttt{OUT\_OF\_SCOPE}; unknown scope returns \texttt{UNKNOWN} with a scope reason;
true scope permits evaluation of the body. A body that would be false does not
turn unknown scope into a known result. The system is answering a scoped
condition, and the applicability question remains unresolved.

Within the body, conjunction and disjunction evaluate all operands in declared
order. This eager behavior retains complete evaluated evidence and exposes errors
or conflicts even if an earlier operand would determine an ordinary Boolean
answer. Missing inputs can appear in a known result, while blocking inputs are
empty when the unknowns do not prevent that result. The distinction is useful:
the bank may know the condition is false while still needing to repair a stale
input feed.

A conditional expression behaves differently. It evaluates the condition and
only the selected branch. An unknown condition returns unknown without evaluating
either branch, even when the branches appear equal. Skipped branches remain type
checked and part of the static dependency closure. Their trace roots are marked
skipped so that the explanation does not invent an evaluation that never occurred.

Defaults add another layer. All exception guards are evaluated first. Two true
guards produce conflict even if their values would agree. An unknown guard can
prevent selection of another true guard because the unknown might become a
competing exception. One true guard with all others false selects its value; all
false guards select the base. Priority is expressed only by explicit nesting with
a reviewed interpretation. Array order is not hidden legal priority.

These rules are not universal truths about law or logic. They are the existing
workbench's specified execution choices. The ensemble may identify a regulation
that needs different semantics. The correct response is to propose and review an
extension or keep the case outside the supported fragment, not to reinterpret
the current operators ad hoc.
% END SOURCE UNIT M07:03

% BEGIN SOURCE UNIT P-03a-rule-contract:02
\subsection{Defaults, branch selection and failure behavior}
\label{merge:P-03a-rule-contract:02}


A default contains a base value and named exceptions. Each exception has a
Boolean guard, a value, an interpretation identifier and source spans. Evaluate
all guards first. Two true guards produce a conflict, even if the proposed values
agree. Otherwise any unknown guard yields an unknown result, because it might
become a competing exception. With exactly one true guard and all others false,
evaluate that exception's value. With all false, evaluate the base. Priority is
expressed only through explicitly reviewed nesting, never row order.

Boolean combinations evaluate all operands in source order. A runtime error or
exception conflict cannot disappear because a different operand is sufficient.
A conditional evaluates its condition and then only the selected branch; an
unknown condition yields unknown, even when both branch values happen to be
equal. Skipped branches are still type checked. These choices eliminate
implementation-dependent short-circuit behavior while retaining precise traces.

The rule's scope is evaluated separately. A false scope yields
\texttt{OUT\_OF\_SCOPE}; an unknown scope yields \texttt{UNKNOWN}. In scope,
Boolean bodies return \texttt{TRUE} or \texttt{FALSE}, other known bodies return
\texttt{VALUE}, and missing evidence returns \texttt{UNKNOWN}. Conflicts and
errors are distinct tagged results. An invalid version time is an error, not a
negative assessment. No exception handler may fall back to false.

The response carries all encountered missing inputs and the conservative subset
blocking the answer. If a known net-asset route establishes the alternative,
the unknown portfolio remains visible but is not a blocker. For an unknown
result, the union of encountered unknown dependencies is reported; the prototype
does not claim it is a mathematically minimal explanation. Traces record evaluated
nodes in postorder, evidence, source spans and explicitly skipped branch roots.
Explanations and diagrams are rendered from that trace.
% END SOURCE UNIT P-03a-rule-contract:02

% BEGIN SOURCE UNIT M07:04
\section{Validation failures are part of observable behavior}
\label{merge:M07:04}


A verifier must compare invalid-input behavior as well as successful decisions.
Malformed JSON, duplicate keys, noncanonical numbers or unpaired Unicode
surrogates can fail at the transport boundary before a semantic result can be
hashed. Valid JSON with an invalid bundle or snapshot can instead produce a
structured semantic error. These outcomes have different evidence and should
not be conflated.

The current validation order is schema, references and duplicates, types, cycles,
time, static-closure conflicts and evaluation. A deterministic order makes error
reports reproducible when an input has several defects. It also prevents a
runtime from accessing an invalid reference before discovering that the bundle
should have been rejected. Fixed error codes and pointers are compared across
backends; dependency-library exception text is not the public contract.

Resource limits are explicit. The present canonical JSON boundary is bounded in
size and depth, and the rule graph has node and expanded-call-depth limits. An
oversized acyclic expression is a resource error; a cycle is a cycle error. An
implementation must not catch either and continue with a truncated rule. The
same principle applies to the new ensemble's graph and report limits.

Tests should combine defects when precedence matters. A snapshot can contain a
conflict and an expired validity interval. A bundle can contain a type error in a
branch that would otherwise be skipped. A default can have two applicable guards
whose values contain errors that should never be evaluated. Isolated operator
tests cannot establish the behavior of these combinations.

A useful specification test states both the expected result and why a tempting
implementation would differ. For example, a short-circuit implementation might
return false before encountering an error in the second operand. The accepted
eager semantics requires the error. This counterexample teaches the contract and
makes the test more durable than an assertion copied from the current function's
output.
% END SOURCE UNIT M07:04

% BEGIN SOURCE UNIT M07:05
\section{Independent oracles and their limits}
\label{merge:M07:05}


An oracle is a source of expected behavior against which an implementation is
compared. Independence has several meanings. A separately written evaluator can
be independent of the production code while relying on the same mistaken
interpretation. A human reference can be independent of the model's answer while
sharing the same incomplete source packet. The verification report must specify
which independence it has actually achieved.

The retained 22-case oracle was frozen before the candidate bundle, which protects
against some forms of adapting expected outputs to the implementation. It was
nevertheless prepared by the same author. The 729-case enumeration gives complete
coverage of six ternary abstract inputs for the selected formula. Neither fact
establishes independent legal interpretation or whole-document coverage.

A stronger reference process begins with independent source reading and scenario
judgment. Reviewers state the expected result, assumptions, supporting passage
and unresolved questions before seeing the candidate output. Disagreements are
adjudicated with reasons or retained as alternative acceptable outcomes. The
reference then has a defined relationship to the evidence rather than being a
single unexplained label.

For engineering checks, a small executable oracle can be valuable when its method
is genuinely different. The second-circular body has a read-once unate structure,
which permits its Boolean-completion check to agree with the specified
three-valued result. The general expression $p\lor\neg p$ shows why that argument
cannot be extended to every RuleIR expression: under the specified semantics,
unknown $p$ yields unknown, while every classical completion yields true.
An oracle based on the wrong generalization would reject correct implementations
or bless incorrect ones.

The oracle itself therefore needs a stated domain and a small set of discriminating
examples. The workbench should retain oracle author, method, source versions,
review status and any shared components. A generated oracle that imports the same
evaluator under a different function name supplies little independent evidence.
A reference that is changed during repair remains useful only when the change is
separately justified and versioned.
% END SOURCE UNIT M07:05

% BEGIN SOURCE UNIT P-03c-storage-review:01
\subsection{The small financial fixture within the complete dry run}
\label{merge:P-03c-storage-review:01}


Section~\ref{sec:dry-run} follows the complete selected transaction profile.
The smaller conformance fixture imports 23EC35 and Annex 1, verifies their hashes and
opens paragraph 3.1 beside the financial rule. Its two wealth definitions are
review issues until the relevant source and data mappings are accepted. An
engineer loads the supplied draft bundle and the synthetic \texttt{f01} case:
portfolio HK\$40 million and qualifying net assets zero. The trace evaluates
both inclusive comparisons and their alternative, returning \texttt{TRUE} for
\texttt{spi.financial}. It also identifies the exact source revision, input
snapshot and engine version.

The reviewer next opens \texttt{f02}, one cent below the portfolio threshold,
with the other route still false. This returns \texttt{FALSE}. In \texttt{f07},
the portfolio is unknown and net assets are HK\$70 million, so the answer is
\texttt{UNKNOWN}. In \texttt{f06}, changing only net assets to HK\$90 million
establishes the financial condition despite the still-missing portfolio. These
cases let a compliance reader inspect the interpretation before any system rule
is exported. A proposed strict comparison is challenged with \texttt{f01}.

Once dependencies and meanings are reviewed, the evidence report is attached to
the bundle hash, reviewers approve it and the release service exports the package.
The production change package includes deterministic Java source and its tested
JAR against the same decision contract. A host adapter supplies the bank's
facts and enforces the released-version and concurrency requirements.
A changed rule creates another hash and invalidates release approval. Historic
decisions still point to the old hash and their original evidence.
% END SOURCE UNIT P-03c-storage-review:01

% BEGIN SOURCE UNIT M07:06
\section{Mutation tests measure what the checks can notice}
\label{merge:M07:06}


A mutation deliberately changes a candidate or implementation in a way that
represents a plausible error. If the verification process rejects it for the
intended reason, the test demonstrates sensitivity to that error. It does not
estimate the frequency of that error in real use, and a high mutation score does
not prove complete correctness.

The second-circular exercise used four meaningful mutations: remove the discount
exception, omit product-type linkage, assume gift classification and ignore scope.
The retained named cases detect all four after compilation. This is more
informative than mutating arbitrary characters because each change corresponds
to a failure a reviewer or programmer could plausibly make.

The ensemble extension should add source and process mutations. Remove a footnote
from one extraction; supply a stale dependency to one member; label a speculative
classification as established; suppress a minority candidate in the report; reset
an issue counter after a revision; or reuse an approval after a source change.
Each mutation targets a distinct assurance boundary. Passing only formula
mutations would leave the new controller largely untested.

Some mutations are equivalent under the selected domain. Changing an expression
may not change any supported outcome, and failure to detect that mutation is not
necessarily a defect. The mutation report must distinguish equivalent mutations,
invalid encodings and live semantic defects. Counting all undetected mutations as
failures, or silently excluding inconvenient ones, would distort the result.

A mutation is killed when a required property changes and the check catches it.
A compiler crash is not always the intended detection: if the task is to test
whether the oracle notices an omitted exception, a syntax-invalid mutant does not
answer the question. Mutation construction should preserve validity where the
research question concerns semantic sensitivity. Invalid-input tests remain a
separate and necessary category.
% END SOURCE UNIT M07:06

% BEGIN SOURCE UNIT P-04-prototype:05
\subsection{Targeted tests and counterexamples}
\label{merge:P-04-prototype:05}


The initial suite should include exact threshold boundaries, small amounts below
them, and cases in which only one financial route holds. It should distinguish
unknown from zero, inconsistent from missing evidence, an absent agreement from
an unretrieved agreement, and client attributes from representative attributes.
These cases test meaning that common ``typical client'' examples fail to expose.

Category tests should change only the relevant category or qualification route.
Ownership tests should change only the allocation agreement or associate status.
Temporal tests should change only the withdrawal sequence, review date or fund
movement. These paired cases help reviewers identify the controlling condition.
For facts designated irrelevant to a narrow rule, changing only those facts should
leave its result unchanged. Such metamorphic checks must name their domain:
adding wealth might be monotone for Equation~\eqref{eq:financial}, but does not
establish monotonicity of every suitability or product-distribution decision.

Mutation tests should deliberately replace inclusive with strict comparisons,
alternatives with conjunctions, category-specific predicates with global ones,
conditional requests with unconditional duties, and named exceptions with generic
fallbacks. Event tests should include duplicated, delayed and out-of-order records.
The test report must identify which defect each case would catch. This borrows
the defect-oriented lesson of CUTECat without claiming that automated case
generation establishes source completeness \citep{cutecat2025}.
% END SOURCE UNIT P-04-prototype:05

% BEGIN SOURCE UNIT M07:07
\section{Symbolic comparison requires a feasible domain}
\label{merge:M07:07}


For supported expressions, symbolic comparison can search for inputs on which two
candidates differ. The domain records admissible values, knownness conditions and
relationships among facts. Before asking whether the outputs differ, the solver
must establish that the domain has at least one model. Otherwise an inconsistent
domain makes every equivalence query vacuously true.

Suppose one domain constraint says that a payment exceeds the fee and another
says that it cannot exceed the same fee. A query finding no differing outcome
under both constraints has taught us nothing about the real candidates. The
workbench returns an inconsistent-domain result. The domain must be repaired or
reviewed before any equivalence claim is made.

A satisfiable difference query returns a witness. The witness must be decoded to
a concrete fact snapshot and replayed in both the reference evaluator and Java.
This replay checks that the symbolic encoding corresponds to executable semantics.
A solver result alone can be wrong relative to the claimed runtime if the
encoding mishandles unknown values, scope or arithmetic. The witness and both
results belong in the report.

An unsatisfiable difference query establishes no differing result within the
checked fragment and domain, subject to the solver and encoding assumptions. It
does not establish that the candidates use the same legal definition, cover the
same source material or behave identically outside the domain. The result label
must therefore include the domain and comparison projection.

The current comparison fragment supports documented Boolean and exact
integer/money reasoning with explicit knownness. It rejects unsupported date,
scale, default and event constructs. The two solver checks share a bounded time
budget, with a separate worker deadline. Unknown, timeout, missing Java replay or
unsupported constructs never become equivalence. The ensemble may schedule
another permitted check, but cannot relabel the incomplete result as a pass.

Counterexample-guided refinement is useful when a witness exposes an imprecise
abstraction. The engineer identifies why the witness is infeasible or why the
encoding omitted a required condition, changes the abstraction with a documented
reason, and reruns the query. The discipline matters: excluding a witness because
it is inconvenient is not refinement. The revised domain must correspond to a
reviewed fact about the application.
% END SOURCE UNIT M07:07

% BEGIN SOURCE UNIT M07:08
\section{Property tests should follow the semantics}
\label{merge:M07:08}


Property-based tests generate many inputs and check a relation that should hold.
They can explore combinations that named examples miss, but the property must
be valid for the specified semantics. Classical Boolean identities cannot all be
transferred to unknown or conflict behavior. Likewise, reordering operands may
preserve a value while changing the required trace or error precedence.

Useful properties include deterministic evaluation for identical canonical
requests, rejection of undeclared facts where the schema requires it, exact
integer arithmetic, preservation of known values under irrelevant evidence
additions, and stable source references under deterministic rendering. Each
property needs a precise precondition. Adding a conflicting record in the static
dependency closure is not an irrelevant addition.

Metamorphic tests compare related inputs when a direct oracle is difficult. If a
known amount increases while all other facts remain fixed, a simple lower-bound
threshold condition should not change from true to false. This monotonicity
applies only to the stated threshold fragment. It need not hold for a default
with an upper-bound exception or for a scope condition that changes at the same
time. The test generator should derive its applicable domain from the property,
not assume the property is universal.

For evidence time, a carefully limited relation is useful: replaying an earlier
\texttt{known\_at} should exclude records first recorded later. Increasing
\texttt{known\_at} does not necessarily preserve the decision, because a correction
or conflict may become visible. A test that demands decision monotonicity in
knowledge time would impose a false property. The intended invariant concerns
which records are admissible, not whether later knowledge always makes the same
answer more certain.

Fuzzing transport boundaries is a different activity. It tests duplicate keys,
size limits, malformed timestamps, unusual Unicode and nesting limits. Its
success criterion is controlled rejection or valid deterministic behavior, not
legal interpretation. These tests protect the integrity of the evidence and
execution chain and should be included in the engineering report under that role.
% END SOURCE UNIT M07:08

% BEGIN SOURCE UNIT P-03d-java-delivery:00
\section{The earlier SPI compiler: a concrete deployment lesson}

The API and commands in the next five subsections belong to the retained
\emph{SPI-Demo1} teaching example under \path{examples/java-dry-run/}. They remain
reproducible, but are not the current v0.1 SDK. The subsequent section on the actual
generated interface supplies the current String/Map API. Keeping both makes the
translation mechanism inspectable without directing an implementer to the wrong
library.


\label{merge:P-03d-java-delivery:00}

\label{sec:java}

Java is the required deployment output. The authoring workbench may use Python
for source processing, review and reference evaluation; a transaction must be
able to call the released Java library without running Python or contacting a
language model. The initial integration assumption is a plain Java 17 JAR.
This keeps the decision code independent of Spring, application servers and
particular rule engines. A bank on another Java version needs a declared target
and the same conformance checks before receiving a compatible build.
% END SOURCE UNIT P-03d-java-delivery:00

% BEGIN SOURCE UNIT P-03d-java-delivery:01
\subsection{SPI-Demo1: the narrower retained demonstration}
\label{merge:P-03d-java-delivery:01}


The delivered \texttt{SPI-Demo1} example contains a source-anchored RuleIR bundle,
synthetic fact snapshots, a separately written Python evaluator, a deterministic
Java emitter, the generated class, two small Java support classes, and a separate
host caller. Its implemented expression subset is deliberately explicit:
Boolean, integer and HK-cent literals and facts; unconditional rule references;
\texttt{all}, \texttt{any}, \texttt{not}; and exact numeric
\texttt{eq}/\texttt{ge}/\texttt{gt} comparisons. The root has a Boolean scope.
The demo compiler rejects dates, arithmetic expressions, general defaults,
conditionals and the full obligation language. The later RuleIR implementation
supports the specified decision operators and its finite duty profile; those
features are outside this earlier demo API. The demonstration response is its own smaller API profile,
not a claim to implement every field and trace node of the full RuleIR service.

The distinction matters for an implementing agent. The example demonstrates an
actual compiler and integration path that can be extended. It does not silently
turn the 35 general RuleIR contract fixtures into executed tests. Those fixtures were the completion target for the full interpreter and
backend and were subsequently executed by the accepted MVP. The separate
SPI example has 32 executed decision scenarios and eleven consent histories, with
its own recorded evidence.
% END SOURCE UNIT P-03d-java-delivery:01

% BEGIN SOURCE UNIT P-03d-java-delivery:02
\subsection{The SPI-Demo1 input and output boundary}
\label{merge:P-03d-java-delivery:02}


The public API accepts an immutable \texttt{Snapshot}, a time being assessed and
a knowledge cutoff. A \texttt{Snapshot} binds a subject identifier to the
declared fact map. A \texttt{Fact} carries a declared scalar type and one of
\texttt{known}, \texttt{unknown} or \texttt{conflict}. A known Boolean is a Java
\texttt{Boolean}; a known monetary or count value is a \texttt{BigInteger}.
Its evidence identifiers, validity interval and recording time are mandatory.
Unknown has a reason and no value. Conflict identifies competing evidence.

The host supplies every declared field, using an explicit unknown where evidence
is unavailable. Omitting a field or supplying the wrong type is a malformed
request, not a legal conclusion. Validity intervals are half-open: evidence
valid until 10:00 does not establish a fact at 10:00. Evidence recorded after
the knowledge cutoff is unavailable to that replay. The API takes canonical UTC
timestamps with six fractional digits; conversion from a Hong Kong business
timestamp belongs in the host adapter and must preserve the intended instant.

\begin{lstlisting}[language=Java,caption={The library API and one dated monetary input.}]
DecisionRuntime.Result evaluate(
    DecisionRuntime.Snapshot snapshot,
    String validAt,
    String knownAt);

// A known value is exact and supported by dated evidence.
Fact portfolio = Fact.known(
    "money_hkd", new BigInteger("4000000000"),
    List.of("valuation.lee", "ownership.lee"),
    "2026-09-01T00:00:00.000000Z",
    "2026-10-01T00:00:00.000000Z",
    "2026-09-01T00:00:00.000000Z");

Fact missingPortfolio = Fact.unknown("money_hkd", "MISSING");
\end{lstlisting}
The first amount is HK\$40 million expressed in cents. A bank adapter reading a
decimal amount should specify the currency, use exact decimal arithmetic and
reject an unsupported fraction of a cent after the adopted conversion. In Java,
\texttt{amount.movePointRight(2).toBigIntegerExact()} can enforce exact cents
after that policy has been applied. Calling \texttt{doubleValue()} first would
discard the property the specification is trying to preserve. For a currency
conversion, the exchange-rate source, date, direction and rounding rule are part
of the evidence mapping.

The result contains a logical status and an operational disposition:
\begin{center}\small
\begin{tabularx}{\textwidth}{P{0.18\textwidth}P{0.34\textwidth}Y}
\toprule Status & Disposition & Host meaning \\\midrule
\texttt{TRUE} & Conditions met & Selected streamlining conditions established; continue the other controls.\\
\texttt{FALSE} & Standard process required & At least one required condition is refuted; do not apply this streamlined route.\\
\texttt{UNKNOWN} & Review required & Evidence can change the conclusion; show blocking inputs.\\
\texttt{CONFLICT} & Review required & Resolve inconsistent evidence; preserve both records.\\
\texttt{OUT\_OF\_SCOPE} & Outside profile & Select another implemented profile or the standard process.\\
\bottomrule
\end{tabularx}
\end{center}
The demo's JSON includes \texttt{profile}, \texttt{authority}, both timestamps,
the bundle and snapshot hashes, a postorder rule trace, missing and blocking
input lists, and the result hash. A rule-trace entry gives its identifier,
status, source-anchor identifiers and blocking inputs. Source anchors resolve
through the bundle included with the release. Known values and their evidence
remain in the preserved snapshot. The current RuleIR implementation provides expression-node traces,
standardized diagnostics and engine identity through the different v0.1 API
documented later in this chapter.

A malformed request throws an explicit Java input exception in this small
library. The production adapter must map that exception, a timeout or an
unavailable release to a technical-failure result and a review/standard-process
path. It must never catch an exception and substitute a successful Boolean.
The demonstration result always carries \texttt{DEMONSTRATION\_ONLY}; production
authority requires a different, approved release record and enforcement by the
host deployment process.
% END SOURCE UNIT P-03d-java-delivery:02

% BEGIN SOURCE UNIT P-03d-java-delivery:03
\subsection{How the compiler emits final Java rather than a suggestion}
\label{merge:P-03d-java-delivery:03}


The emitter first validates the JSON schema, references, identifiers, types and
acyclic rule graph. It then assigns an internal Java method name to each rule
and recursively emits its expression. Identifiers are validated data, never
pieces of executable text supplied by a language model. The generated monetary
comparison is represented by actual Java calls such as:
\begin{lstlisting}[language=Java,caption={The generated financial method's essential body. The complete class is built and tested in the supplied example.}]
return c.rule("spi.financial", sourceIds, () ->
    any(
        compare("ge", c.fact("portfolio"),
            Value.of(new BigInteger("4000000000"))),
        compare("ge", c.fact("net_assets_ex_home"),
            Value.of(new BigInteger("8000000000")))
    ));
\end{lstlisting}
\texttt{compare} returns a known Boolean when both operands are known, or an
unknown with its blocking input names. \texttt{any} implements the truth table
in Equation~\eqref{eq:or}. Java evaluates its arguments before calling the
method, so both branches have the specified eager evaluation and evidence
collection. Replacing this with Java's short-circuit \texttt{||} would change
the trace and missing-input behavior even where the final Boolean agreed.
The context memoizes referenced rules and records their results in a stable
order. A precheck detects conflicting facts in this profile's dependency set
before normal expression evaluation.

The production emitter should implement one mapping per supported operator,
with an explicit error for every unsupported construct. It must not ask an LLM
to write new Java on each regulatory change. The LLM can propose the typed
specification; after review, a deterministic compiler performs the mechanical
translation. This concentrates semantic review on a small language and makes
the generation step reproducible.

For the money comparison, the preservation argument is short. The IR represents
both amounts as mathematical integers of HK cents. The Java boundary constructs
the same integers as \texttt{BigInteger} values. The emitted comparison returns
true exactly when \texttt{compareTo} is nonnegative, which expresses the same
inclusive ordering. The argument depends on the input mapping preserving those
integers. It would fail if the adapter truncated decimals, mixed currencies or
overflowed a narrower integer before constructing the value. For a complete
compiler, analogous local arguments must be composed for every operator,
unknown/error case, rule reference and trace requirement. The delivered tests
check examples of preservation; they are not that universal proof.
% END SOURCE UNIT P-03d-java-delivery:03

% BEGIN SOURCE UNIT P-03d-java-delivery:04
\subsection{Reproduce the build and see the host result}
\label{merge:P-03d-java-delivery:04}


The repository retains a project-local Temurin JDK 17 for this check, verified
against the official archive checksum. It does not change the machine's default
Java installation. From the repository root, the single build command is:

\begin{minipage}{\textwidth}
\begin{lstlisting}
python3 examples/java-dry-run/build_and_verify.py \
  --jdk .localresources/java-toolchain/jdk-17.0.20.1+1
\end{lstlisting}
\end{minipage}

The Python authoring check uses the already available JSON Schema package. The
generated Java library itself has no third-party runtime dependencies. The build
uses \texttt{javac --release 17 -Xlint:all -Werror}; creates the JAR; compiles
the separate fixture and host callers against that JAR; compares complete
reference and Java results; and runs the three mutations. The package has fixed
ZIP timestamps to make repeated builds with the same toolchain reproducible.
The verification record stores compiler identity, source/compiler/JAR digests,
the executed command, results and limitations.

The delivered files below are operational copies of the design, not additional
reading needed to understand why it works:
\begin{center}\small
\begin{tabularx}{\textwidth}{P{0.45\textwidth}Y}
\toprule File under \path{examples/java-dry-run/} & Purpose \\\midrule
\path{spec/spi-control.bundle.json} & Source-linked typed interpretation for the chosen profile.\\
\path{spec/decision-cases.json}; \path{spec/consent-cases.json} & Independent expected statuses and synthetic evidence/history inputs.\\
\path{reference.py}; \path{compile_java.py} & Separate reference evaluation and deterministic Java generation.\\
\path{generated/GeneratedSpi.java} & Final generated decision class.\\
\path{src/main/java/hk/legalmath/spi/} & Evidence, result, three-valued arithmetic/logic support and consent replay.\\
\path{build/spi-controls-demo.jar} & Compiled library callable by a Java host.\\
\path{generated/BankHostExample.java} & Complete synthetic host caller using the public API.\\
\path{build/verification.json}; \path{build/host-output.txt} & Executed validation and host results.\\
\bottomrule
\end{tabularx}
\end{center}

The separate host prints these dispositions for the synthetic history:
\begin{lstlisting}
BEFORE_WITHDRAWAL STREAMLINING_CONDITIONS_MET
AFTER_WITHDRAWAL STANDARD_PROCESS_REQUIRED
RESULT_AUTHORITY DEMONSTRATION_ONLY
\end{lstlisting}
This is the concrete answer to how final code reaches Java: the bank depends on
the released library, constructs the declared evidence snapshot through its
adapter, calls the generated class and handles its structured result. The
compiler, public circulars and drafting tools are build-time/review-time
components. They need not be deployed into the bank's order-processing process.
% END SOURCE UNIT P-03d-java-delivery:04

% BEGIN SOURCE UNIT M07:09
\section{The actual generated Java interface}
\label{merge:M07:09}


The current emitter validates the bundle, canonicalizes it and embeds its exact
bytes in a generated class. The class name is
\texttt{hk.legalmath.Policy\_} followed by the first twenty hexadecimal digits of
the bundle hash. The class exposes the full bundle hash and delegates evaluation
to the Java semantic runtime. This is a generated, immutable policy binding; it
is not the older proposal's imagined record-based API or a proof that each legal
rule has been translated to hand-optimized Java control flow.

The two current entry points are:
\begin{lstlisting}[language=Java,caption={Implemented generated-policy API.}]
public static String evaluate(
    String snapshotJson,
    String ruleId,
    String validAt,
    String knownAt,
    String mode);

public static Map<String,Object> evaluate(
    Map<String,Object> snapshot,
    String ruleId,
    String validAt,
    String knownAt,
    String mode);
\end{lstlisting}

The string overload returns strict JSON. The map overload returns a deeply
immutable result. The generated policy and Java runtime operate without Python,
a language model, a solver or network retrieval. This property is important for
a bank integration: the interpretation workbench prepares and verifies a package,
while production decision execution can remain local and deterministic.

A concrete host compiles against the actual generated class name from the build
manifest. It supplies a validated snapshot, exact rule identifier, both assessment
times and the requested mode. The mode does not grant authority. A caller cannot
turn a draft into an approved control by passing the word ``production.'' The
release service and host must establish approval and applicability separately.

The conceptual integration looks as follows. The symbolic class name in the
listing is replaced by the actual generated name when the integration fixture is
built; the release lookup and transaction methods are host responsibilities.
\begin{lstlisting}[language=Java,caption={Host responsibilities around deterministic evaluation.}]
ApprovedRelease release = releases.resolve(profile, validAt, knownAt);
release.verifyAuthorityAndDigests();
ConsistentSnapshot s = adapter.readSnapshotWithVersions(subject, order);
String resultJson = Policy_HASH.evaluate(
    s.strictJson(), release.ruleId(), validAt, knownAt, "production");
Decision decision = decoder.decodeStrictly(resultJson);
HostDisposition disposition = hostPolicy.map(decision);
transactions.commitIfVersionsUnchanged(
    s.versions(), order.id(), order.payloadHash(),
    release.identity(), decision, disposition);
\end{lstlisting}

\texttt{ApprovedRelease}, \texttt{ConsistentSnapshot} and the host transaction
interfaces in this illustration are proposed integration concepts, not classes
already supplied by the generated SDK. Their responsibilities must be implemented
in the bank's environment. The example makes the boundary explicit so an agent
does not mistake illustrative types for an existing dependency.

The host must handle every tagged result. A true selected prohibition condition
may require stopping the relevant offer; false permits only progression to other
controls under the host policy. Unknown, conflict, out-of-scope and error each
have explicit dispositions. No result string should be cast to a Boolean by a
language's truthiness rules. The decoder rejects unknown status tags so that a
future runtime extension cannot silently take an old default branch.
% END SOURCE UNIT M07:09

% BEGIN SOURCE UNIT M07:10
\section{Traces are part of the delivered meaning}
\label{merge:M07:10}


A decision without an explanation is difficult to investigate. The runtime trace
records node identifiers, operators, children, type, value or status, evidence
references and source spans. Shared rule evaluations are memoized, and skipped
branches have explicit markers. The trace is deterministic under the request and
engine version, making it suitable for replay and comparison.

Trace fidelity is a separate verification obligation. Two programs may return
true while one omits the evidence for the exception or incorrectly claims to have
evaluated a skipped branch. A regulator-facing explanation assembled from that
trace could be misleading even though the Boolean result matches. The comparison
therefore includes trace structure and evidence, not just the top-level status.

Natural-language explanations are generated from the verified trace and approved
source mappings. A language model may help render a readable summary, but it
cannot alter the computed result or invent a justification. If a summary mentions
a source that the trace and bundle do not contain, validation should reject or
flag it. The authoritative explanation remains reconstructible without the
summary model.

The current source-span policy associates trace nodes with the owning rule's
explicit citations and retains exception citations in the bundle. That is a
coarser mapping than a separate legal proof for every arithmetic operation. The
report should describe the granularity honestly. A future refinement may attach
more precise spans, but a dense set of citations should not be mistaken for a
formal demonstration that each mapping is legally correct.

Execution hashes identify the canonical request and result under a particular
backend version. They help detect altered bytes and reproduce behavior. They do
not authenticate a reviewer or prove that a document came from the regulator.
Authentication, source acquisition and approval are separate controls. Confusing
these roles is a common source of overclaiming in systems described as auditable.
% END SOURCE UNIT M07:10

% BEGIN SOURCE UNIT M07:11
\section{Event histories add obligations that a snapshot cannot express}
\label{merge:M07:11}


Some requirements concern sequences: obtain consent before an action, honor a
withdrawal, provide a document within a period, or preserve a breach after late
performance. A current Boolean snapshot cannot express all of these relationships.
The existing event component replays finite attributed histories using a separately
versioned Java entry point, \texttt{EventReplay.replay(Map<String,Object>)}.

A stream header binds subject, category, actor, action, profile and inception.
Events belong to that header; an API path cannot redirect their meaning to a
different subject. Duplicate identical events are idempotent, while conflicting
payloads or sequence reuse are errors. Ordering authority is explicit. If event
order is unresolved, the system must not invent a sequence based solely on arrival.

Consent illustrates why completeness matters. Absence of a grant in an incomplete
history does not establish that no grant occurred. A completeness watermark or a
reviewed initial state supplies the relevant observation boundary. A watermark
cannot cover a time later than the recording instant of its assertion. A
withdrawal ends the applicable consent generation, and a later regrant uses a new
evidence identity rather than resurrecting a previously withdrawn record.

For an achievement obligation, activation opens a duty with a deadline. The
selected contract uses the interval from activation up to, but excluding, the
deadline for timely performance. Performance exactly at the excluded endpoint is
therefore late under that contract. This is an engineering convention that must
match the reviewed legal timing rule; a different legal rule requires a different
specified boundary, not a casual adjustment in the host.

Late performance does not erase the historical breach. The system can record both
that a duty was breached and that it was later fulfilled. Similarly, an observed
action is retained even if an advisory control would have prohibited it. Deleting
inconvenient observations would make replay describe an idealized compliant world
rather than the bank's actual history.

Each opened duty binds the source bundle that created it. A later amendment does
not automatically migrate existing duties. Migration may be appropriate, but it
requires an explicit legal and operational policy about which obligations change,
which dates apply and how historical findings are retained. The current MVP has
no automatic migration implementation. The ensemble can identify the question
and prepare alternatives; it cannot assume the answer from the new publication
date alone.
% END SOURCE UNIT M07:11

% BEGIN SOURCE UNIT P-03-product:04
\subsection{Obligations and event histories}
\label{merge:P-03-product:04}


The decision fragment cannot represent every continuing control. For consent,
let $s_i$ describe the recorded arrangement immediately before event $e_i$:
its client, selected categories, agreed threshold, monitoring mode, relevant
reviews and withdrawal status. Each event includes its timestamps and ordering
evidence. It changes the recorded state through a
versioned transition function,
\begin{equation}
 s_{i+1}=T_v(s_i,e_i).
 \label{eq:transition}
\end{equation}
Here $i$ orders processed events; it does not assume equal physical time
steps. Version $v$ identifies the adopted rules and ordering convention. A
withdrawal event can change the availability of the streamlined arrangement for
subsequent decisions. A new-funds event can trigger the designated-account
checks. An annual-review event creates work that must later be evidenced as
completed. The exact consequences require review of the cited guidance and the
bank's operating procedure \citep{sfcspiannex1,sfcspiannex2}.

An obligation record should distinguish its creation, due condition, performance,
breach and any later remediation. If a transaction occurs after a withdrawal, the
audit view must retain the transaction and identify the issue; it cannot erase
the event because the authorization view would have blocked it. This follows the
useful separation of normative permission from actual occurrence in eFLINT
\citep{eflint2026}.

Time requires two histories. Legal validity time records when a source or adopted
rule applies. System-recorded time records when the bank received and deployed
it. These support different questions: what should have applied to an event, and
what version did the system actually use? The same distinction applies to facts
received late or corrected later. A reconstructed historical assessment may use
new evidence, but it must not overwrite the original decision and its inputs.

Event ordering is a material assumption. The prototype should store source
timestamps, receipt timestamps, sequence identifiers, time zone and any uncertainty
in ordering. A synthetic trace with a consent withdrawal and execution at the
same timestamp is unresolved unless the available ordering or an adopted
convention settles it. Calendar computations should identify whether they use
anniversaries, elapsed days or another reviewed operation, drawing on the date
semantics literature \citep{dates2024}. Neither the operating system clock nor the
order in which messages happen to arrive establishes legal priority.
% END SOURCE UNIT P-03-product:04

% BEGIN SOURCE UNIT P-03b-event-contract:00
\subsection{The finite event profile}
\label{merge:P-03b-event-contract:00}


The first event implementation supports consent and non-preemptive achievement
duties. Events have immutable identifiers, stream identifiers, occurrence and
recorded timestamps, and an authoritative sequence number. Timestamps are UTC
with microsecond precision; legal date conversion explicitly uses
\texttt{Asia/Hong\_Kong}. Equal occurrence times can use sequence order only
when the connector declares that ordering authority. Otherwise the result is
\texttt{E\_ORDER\_UNRESOLVED}. A duplicate event with identical bytes is a
no-op; reuse of its identifier for different content is an error.

Consent begins as not established only when the stream is complete from its
declared inception, or is complete after a reviewed initial snapshot. An incomplete export
instead yields an unknown history. A grant creates an active consent generation;
withdrawal makes it withdrawn; a new grant needs new evidence and creates a new
generation. Observed transactions remain in the audit history even when the
current consent state would prevent streamlined use.

An obligation records actor, action, subject, source version, activation and a
deadline or explicit absence of a deadline. Let $a$ be activation and $d$ the
exclusive deadline. Under this proposed profile, a matching performance at time
$p$ is timely when
\begin{equation}
 a\leq p<d.
 \label{eq:timely}
\end{equation}
This is the chosen boundary convention for the profile, not an assertion that
every SFC deadline has the same wording. A local due date is converted to the
next local midnight. An anniversary requires an explicit leap-day policy; a
reviewed choice of February 28, March 1 or rejection replaces an implicit library
default. Business-day calendars are deferred.

Absence of a timely record establishes breach only when a watermark shows that
the observation window through $d$ is complete. The wall clock cannot create
that evidence. Without completeness the obligation stays open. Timely performance
produces \texttt{SATISFIED\_ON\_TIME}; a complete window without it produces
\texttt{BREACHED}. A subsequent matching performance produces
\texttt{PERFORMED\_LATE} and retains the breach. A deadline-free request can be
performed but does not become overdue through an invented number of days.

For example, consider a synthetic review activated at midnight UTC on September
1 and due at midnight on September 2. A record one microsecond before the latter
instant is timely; a record exactly at it is late. The fixtures contain both.
These dates illustrate the runtime convention rather than an SFC-prescribed
review deadline. If a late-arriving record later proves timely performance,
replay produces a corrected assessment with a new knowledge cutoff; the original
assessment remains visible. Maintenance duties, pre-emptive satisfaction,
waivers and legal compensation require future profiles and are rejected now.
% END SOURCE UNIT P-03b-event-contract:00

% BEGIN SOURCE UNIT M07:12
\section{Build, verification and approval form an acyclic chain}
\label{merge:M07:12}


A build manifest records the bundle, generated source, runtime source, toolchain,
Java target, dependencies and JAR digest. A verification report then names the
build and exact JAR it tested. A release manifest names that build and verification
report, together with applicability and effective-time information. Approval
binds the final release manifest and bundle.

The direction matters. Embedding a manifest that contains the JAR's own digest
inside that JAR creates a circular identity problem. The current design keeps
these records external and acyclic. A reviewer approves the exact tested bytes
through the final manifest, rather than approving a mutable directory or a
human-readable version label.

Candidate compilation occurs before legal approval, because reviewers need a
concrete package and verification evidence to inspect. Compilation does not grant
deployment authority. Export into a controlled release route occurs only after
the required approvals and checks. This distinction lets engineers do useful
reversible work without repeatedly asking for permission, while preserving the
human decision at the point where it matters.

Reproducible builds strengthen the identity chain. Under the same declared
toolchain and inputs, build in separate directories with fixed archive-entry
timestamps and compare emitted source and binary hashes. A mismatch should be
investigated and recorded. Reproducibility does not prove semantic correctness,
but it reduces uncertainty about which bytes correspond to the reviewed source.

The local MVP uses synthetic identities and unsigned historical archives to
exercise workflow. A bank deployment needs enterprise identity, role separation,
key management and independently verified release signatures or equivalent
controls. A hash proves equality of bytes under its cryptographic assumption;
it does not prove who authorized those bytes. The manuscript's implementation
plan keeps that integration distinct from the local demonstration.
% END SOURCE UNIT M07:12

% BEGIN SOURCE UNIT M07:13
\section{A correct decision can still race with the transaction}
\label{merge:M07:13}


Suppose a control depends on consent and exposure. The host reads both, evaluates
an order and then commits it. Between the read and commit, consent may be
withdrawn or another order may consume the available limit. The decision was
correct for its snapshot but is no longer valid for the transaction being
committed. This is a concurrency problem outside the pure rule evaluator.

The host should read a consistent snapshot with version identifiers, evaluate
exposure including the proposed order, and atomically validate the versions while
reserving the exposure and recording the decision. If the versions changed, it
rereads and reevaluates under a bounded transaction retry policy. That policy is
separate from interpretation-resolution retries and must not share an ambiguous
counter called simply ``attempts.''

Idempotency protects against duplicate submissions. Repeating the same order
identifier and payload should recover the same committed result or reservation.
Using the same identifier with a different payload is a conflict. The identity
must include the legally relevant order content, not merely a user-interface
request number. Otherwise a retry can accidentally authorize a changed transaction.

Release selection has a similar race. The host records the release identity used
for the decision and verifies that the applicable release policy permits it at
commit. Whether an in-flight transaction can finish under a prior release is a
bank policy requiring explicit treatment. The rule engine cannot resolve that
question by consulting its process clock or loading whatever JAR is newest.

A rollback also has legal meaning. Restoring a previous software build can repair
an implementation defect if that build still expresses the applicable rule.
Restoring a superseded legal regime is different. The incident procedure must
separate software rollback, operational suspension and legal-rule selection.
Historical packages remain available for replay even when they cannot be used for
new decisions.
% END SOURCE UNIT M07:13

% BEGIN SOURCE UNIT P-03d-java-delivery:05
\subsection{The bank adapter and the exposure race}
\label{merge:P-03d-java-delivery:05}


The adapter has four responsibilities. It selects the approved rule release for
the legal entity, business profile and relevant time; resolves identities,
categories and evidence into the declared inputs; obtains a consistent consent
and exposure snapshot; and records the result alongside the transaction's other
controls. The generated library should have no credentials for editing these
source systems. A host record should include transaction identifier, release and
snapshot hashes, category and consent generation, data-version identifiers,
decision timestamps and the final disposition.

Concurrency is a material part of correctness. Suppose existing gross exposure
is HK\$8 million and two orders each add HK\$1.5 million against a HK\$10 million
limit. If each reads HK\$8 million, both independently pass at HK\$9.5 million,
but together they produce HK\$11 million. A perfectly correct generated formula
does not prevent that race. The host must include reserved exposure and commit
the new reservation under a transaction or optimistic version check. If the
exposure or consent version changed after evaluation, it reloads and evaluates
again. A withdrawal must invalidate a stale consent version before a new order
can consume the decision. Market-exposure policy, release timing and retry limits
must be agreed with the bank's order and risk owners.

An implementable host sequence is:
\begin{enumerate}
\item Resolve an idempotency key from the order request. A retry with the same
key and unchanged payload returns its recorded disposition; a changed payload
under that key conflicts.
\item Read the approved release and the subject/category consent generation,
holdings, leverage and outstanding reservations with their version identifiers.
Build the dated snapshot and calculate exposure including this order.
\item Evaluate the generated Java library and the host's remaining controls.
If this profile is unresolved or refuted, do not consume its streamlined route.
\item Atomically validate the read versions, reserve the permitted exposure and
append the decision record. On version conflict, repeat from the read step with
a bounded retry policy; otherwise route to review.
\item Retain the committed record for replay. A later correction creates another
assessment linked to the old one; it does not overwrite what was known earlier.
\end{enumerate}
The standalone JAR demonstrates the decision step. It does not simulate this
bank transaction protocol or certify concurrency in an unspecified host system.
Those integration tests are required before production use.
% END SOURCE UNIT P-03d-java-delivery:05

% BEGIN SOURCE UNIT M07:14
\section{A verification dossier that supports a release decision}
\label{merge:M07:14}


The final technical dossier links the accepted interpretation to its implementation
through several kinds of evidence. Named scenarios demonstrate reviewed boundary
behavior. Exhaustive finite checks cover declared small domains. Property and
mutation tests examine broader classes of implementation mistakes. Symbolic
comparison supplies domain-qualified results and replayed witnesses. Differential
execution compares Python and Java semantics. Build records identify exact bytes.
Host tests exercise concurrency, idempotency and release selection.

These methods overlap but are not interchangeable. A compiler proof could cover
an unbounded class of accepted expressions under stated semantics, while still
assuming that the factual adapter and source interpretation are correct. A test
suite can catch a concrete adapter error that lies outside that proof. An
independent legal review can reject a perfectly verified expression because it
misreads an exception. The dossier should state the role of each result so that
one strong component does not hide a missing one.

For the current project, the retained engineering evidence includes the local
MVP test suite and the separate second-circular exercise. The new ensemble is not
yet an implemented live-model product. The appropriate next engineering step is
to implement its controller and contracts, then run deterministic acceptance
scenarios, then evaluate interpretation performance with independent reference
judgments. Declaring the entire product ready because the Java path works would
skip the central problem this book is designed to address.

A release reviewer should be able to answer four concrete questions from the
package: what rule is being enforced; what evidence makes its inputs meaningful;
what demonstrates that these exact bytes preserve the rule; and what happens when
an assumption or input is uncertain. If any answer requires guessing from a model
confidence score, the dossier is incomplete. Chapter~\ref{ch:implementation} turns
these requirements into service records, operations and ordered implementation
tasks.
% END SOURCE UNIT M07:14

